% =====================================================================
% Note on: Why We Don't Discount at the "Risk-Free Rate" (anonymous)
% =====================================================================
\documentclass[11pt,a4paper]{article}

\newcommand{\papertag}{anon\_discounting\_textbooks}

% =====================================================================
% Shared preamble for paper notes
% Include with:  % =====================================================================
% Shared preamble for paper notes
% Include with:  % =====================================================================
% Shared preamble for paper notes
% Include with:  \input{../_common/preamble}
% =====================================================================

% --- Core packages ---
\usepackage[utf8]{inputenc}
\usepackage[T1]{fontenc}
\usepackage{lmodern}
\usepackage[margin=2.5cm]{geometry}
\usepackage{amsmath,amssymb,amsthm,mathtools}
\usepackage{bm}
\usepackage{booktabs}
\usepackage{array}
\usepackage{enumitem}
\usepackage{hyperref}
\usepackage{xcolor}
\usepackage{listings}
\usepackage{fancyhdr}
\usepackage{titlesec}

% --- Hyperref styling ---
\hypersetup{
  colorlinks=true,
  linkcolor=blue!50!black,
  urlcolor=blue!50!black,
  citecolor=blue!50!black
}

% --- Theorem-like environments ---
\theoremstyle{definition}
\newtheorem{defn}{Definition}
\newtheorem{remark}{Remark}
\newtheorem{example}{Example}

\theoremstyle{plain}
\newtheorem{prop}{Proposition}
\newtheorem{lemma}{Lemma}

% --- Coloured boxes for the structured sections ---
% Validation block, commentary, TODOs use coloured frames so the eye
% can find them when flipping through the note.
\usepackage[most]{tcolorbox}

\newtcolorbox{commentary}{
  colback=yellow!5,
  colframe=yellow!50!black,
  title=Commentary,
  fonttitle=\bfseries,
  breakable
}

\newtcolorbox{validation}{
  colback=green!3,
  colframe=green!50!black,
  title=Validation,
  fonttitle=\bfseries,
  breakable
}

\newtcolorbox{todobox}{
  colback=red!3,
  colframe=red!50!black,
  title=TODO / Open questions,
  fonttitle=\bfseries,
  breakable
}

% --- Code listing style for Python snippets in validation blocks ---
\lstdefinestyle{py}{
  language=Python,
  basicstyle=\ttfamily\footnotesize,
  keywordstyle=\color{blue!60!black},
  commentstyle=\color{gray},
  stringstyle=\color{red!60!black},
  showstringspaces=false,
  breaklines=true,
  frame=single,
  framesep=4pt,
  rulecolor=\color{gray!30}
}

% --- Page header: shows paper tag so a printed note is identifiable ---
\pagestyle{fancy}
\fancyhf{}
\fancyhead[L]{\small\textit{\papertag}}
\fancyhead[R]{\small\thepage}
\renewcommand{\headrulewidth}{0.4pt}

% --- Section spacing: tighter than article default ---
\titlespacing*{\section}{0pt}{1.5ex plus 0.5ex minus 0.2ex}{1ex plus 0.2ex}
\titlespacing*{\subsection}{0pt}{1.2ex plus 0.3ex minus 0.2ex}{0.6ex plus 0.2ex}

% =====================================================================

% --- Core packages ---
\usepackage[utf8]{inputenc}
\usepackage[T1]{fontenc}
\usepackage{lmodern}
\usepackage[margin=2.5cm]{geometry}
\usepackage{amsmath,amssymb,amsthm,mathtools}
\usepackage{bm}
\usepackage{booktabs}
\usepackage{array}
\usepackage{enumitem}
\usepackage{hyperref}
\usepackage{xcolor}
\usepackage{listings}
\usepackage{fancyhdr}
\usepackage{titlesec}

% --- Hyperref styling ---
\hypersetup{
  colorlinks=true,
  linkcolor=blue!50!black,
  urlcolor=blue!50!black,
  citecolor=blue!50!black
}

% --- Theorem-like environments ---
\theoremstyle{definition}
\newtheorem{defn}{Definition}
\newtheorem{remark}{Remark}
\newtheorem{example}{Example}

\theoremstyle{plain}
\newtheorem{prop}{Proposition}
\newtheorem{lemma}{Lemma}

% --- Coloured boxes for the structured sections ---
% Validation block, commentary, TODOs use coloured frames so the eye
% can find them when flipping through the note.
\usepackage[most]{tcolorbox}

\newtcolorbox{commentary}{
  colback=yellow!5,
  colframe=yellow!50!black,
  title=Commentary,
  fonttitle=\bfseries,
  breakable
}

\newtcolorbox{validation}{
  colback=green!3,
  colframe=green!50!black,
  title=Validation,
  fonttitle=\bfseries,
  breakable
}

\newtcolorbox{todobox}{
  colback=red!3,
  colframe=red!50!black,
  title=TODO / Open questions,
  fonttitle=\bfseries,
  breakable
}

% --- Code listing style for Python snippets in validation blocks ---
\lstdefinestyle{py}{
  language=Python,
  basicstyle=\ttfamily\footnotesize,
  keywordstyle=\color{blue!60!black},
  commentstyle=\color{gray},
  stringstyle=\color{red!60!black},
  showstringspaces=false,
  breaklines=true,
  frame=single,
  framesep=4pt,
  rulecolor=\color{gray!30}
}

% --- Page header: shows paper tag so a printed note is identifiable ---
\pagestyle{fancy}
\fancyhf{}
\fancyhead[L]{\small\textit{\papertag}}
\fancyhead[R]{\small\thepage}
\renewcommand{\headrulewidth}{0.4pt}

% --- Section spacing: tighter than article default ---
\titlespacing*{\section}{0pt}{1.5ex plus 0.5ex minus 0.2ex}{1ex plus 0.2ex}
\titlespacing*{\subsection}{0pt}{1.2ex plus 0.3ex minus 0.2ex}{0.6ex plus 0.2ex}

% =====================================================================

% --- Core packages ---
\usepackage[utf8]{inputenc}
\usepackage[T1]{fontenc}
\usepackage{lmodern}
\usepackage[margin=2.5cm]{geometry}
\usepackage{amsmath,amssymb,amsthm,mathtools}
\usepackage{bm}
\usepackage{booktabs}
\usepackage{array}
\usepackage{enumitem}
\usepackage{hyperref}
\usepackage{xcolor}
\usepackage{listings}
\usepackage{fancyhdr}
\usepackage{titlesec}

% --- Hyperref styling ---
\hypersetup{
  colorlinks=true,
  linkcolor=blue!50!black,
  urlcolor=blue!50!black,
  citecolor=blue!50!black
}

% --- Theorem-like environments ---
\theoremstyle{definition}
\newtheorem{defn}{Definition}
\newtheorem{remark}{Remark}
\newtheorem{example}{Example}

\theoremstyle{plain}
\newtheorem{prop}{Proposition}
\newtheorem{lemma}{Lemma}

% --- Coloured boxes for the structured sections ---
% Validation block, commentary, TODOs use coloured frames so the eye
% can find them when flipping through the note.
\usepackage[most]{tcolorbox}

\newtcolorbox{commentary}{
  colback=yellow!5,
  colframe=yellow!50!black,
  title=Commentary,
  fonttitle=\bfseries,
  breakable
}

\newtcolorbox{validation}{
  colback=green!3,
  colframe=green!50!black,
  title=Validation,
  fonttitle=\bfseries,
  breakable
}

\newtcolorbox{todobox}{
  colback=red!3,
  colframe=red!50!black,
  title=TODO / Open questions,
  fonttitle=\bfseries,
  breakable
}

% --- Code listing style for Python snippets in validation blocks ---
\lstdefinestyle{py}{
  language=Python,
  basicstyle=\ttfamily\footnotesize,
  keywordstyle=\color{blue!60!black},
  commentstyle=\color{gray},
  stringstyle=\color{red!60!black},
  showstringspaces=false,
  breaklines=true,
  frame=single,
  framesep=4pt,
  rulecolor=\color{gray!30}
}

% --- Page header: shows paper tag so a printed note is identifiable ---
\pagestyle{fancy}
\fancyhf{}
\fancyhead[L]{\small\textit{\papertag}}
\fancyhead[R]{\small\thepage}
\renewcommand{\headrulewidth}{0.4pt}

% --- Section spacing: tighter than article default ---
\titlespacing*{\section}{0pt}{1.5ex plus 0.5ex minus 0.2ex}{1ex plus 0.2ex}
\titlespacing*{\subsection}{0pt}{1.2ex plus 0.3ex minus 0.2ex}{0.6ex plus 0.2ex}

% =====================================================================
% House notation: shared symbols used across all paper notes.
% Each note imports this and may shadow individual macros if a paper
% uses a clashing convention — but the *default* meaning is fixed here.
%
% Include with:  % =====================================================================
% House notation: shared symbols used across all paper notes.
% Each note imports this and may shadow individual macros if a paper
% uses a clashing convention — but the *default* meaning is fixed here.
%
% Include with:  % =====================================================================
% House notation: shared symbols used across all paper notes.
% Each note imports this and may shadow individual macros if a paper
% uses a clashing convention — but the *default* meaning is fixed here.
%
% Include with:  \input{../_common/notation}
% =====================================================================

% --- Measures and expectations ---
\newcommand{\Q}{\mathbb{Q}}              % risk-neutral measure
\newcommand{\Qt}{\mathbb{Q}^{T}}         % T-forward measure
\newcommand{\E}{\mathbb{E}}              % expectation
\newcommand{\Et}[1]{\E^{#1}}             % expectation under measure #1
\newcommand{\Ftime}[1]{\mathcal{F}_{#1}} % filtration at #1

% --- Discount and forward objects ---
\newcommand{\Disc}[2]{D_{#1,#2}}         % discount factor D_{t,T}
\newcommand{\Bond}[2]{P_{#1}(#2)}        % bond price P_t(y)
\newcommand{\Ann}[1]{A_{#1}}             % annuity / level
\newcommand{\Numer}[1]{N_{#1}}           % numeraire
\newcommand{\Fwd}[2]{F_{#1,#2}}          % forward price F_{t,T}
\newcommand{\Fwdpx}[2]{\mathrm{Fwd}_{#1}(#2)} % verbose forward-price F_t(T)

% --- Rates ---
\newcommand{\IRR}{\mathrm{IRR}}          % internal rate of return
\newcommand{\Yld}{y}                     % bond yield variable
\newcommand{\Strike}{K}                  % strike (price-space)
\newcommand{\Lock}{L}                    % locked yield / rate
\newcommand{\Repo}{r^{\mathrm{repo}}}    % repo rate
\newcommand{\Fund}{r^{\mathrm{fun}}}     % unsecured funding rate
\newcommand{\OIS}{r^{\mathrm{ois}}}      % OIS rate
\newcommand{\Coll}{r^{\mathrm{c}}}       % collateral rate
\newcommand{\Rcmt}{R^{\mathrm{cmt}}}     % constant-maturity-treasury rate
\newcommand{\Rswp}{R^{\mathrm{swp}}}     % swap rate (used for risky variant)
\newcommand{\Rec}{\mathrm{Rec}}          % recovery rate
\newcommand{\NPV}{\mathrm{NPV}}          % present-value functional
\newcommand{\PVone}{\mathrm{PV01}}       % risky annuity
\newcommand{\Loss}{\mathrm{Loss}}        % default-leg integral
\newcommand{\Sasw}{S^{\mathrm{ASW}}}     % asset-swap spread
\newcommand{\Scds}{S^{\mathrm{CDS}}}     % CDS spread
\newcommand{\Basis}{\mathrm{Basis}}      % CDS-bond basis

% --- Sensitivities ---
\newcommand{\Diff}[2]{D_{#1}^{#2}}       % Euler's notation: D_x^k[f]
\newcommand{\Riskf}{\mathrm{RiskFactor}} % bond risk factor (= -DV01*1e4)
\newcommand{\DVone}{\mathrm{DV01}}

% --- Indicator / misc ---
\newcommand{\Ind}{\mathbf{1}}
\newcommand{\given}{\,|\,}
\newcommand{\dd}{\,\mathrm{d}}

% --- Greeks-style ---
\newcommand{\Gam}{\Gamma}
\newcommand{\Vega}{\mathcal{V}}

% =====================================================================

% --- Measures and expectations ---
\newcommand{\Q}{\mathbb{Q}}              % risk-neutral measure
\newcommand{\Qt}{\mathbb{Q}^{T}}         % T-forward measure
\newcommand{\E}{\mathbb{E}}              % expectation
\newcommand{\Et}[1]{\E^{#1}}             % expectation under measure #1
\newcommand{\Ftime}[1]{\mathcal{F}_{#1}} % filtration at #1

% --- Discount and forward objects ---
\newcommand{\Disc}[2]{D_{#1,#2}}         % discount factor D_{t,T}
\newcommand{\Bond}[2]{P_{#1}(#2)}        % bond price P_t(y)
\newcommand{\Ann}[1]{A_{#1}}             % annuity / level
\newcommand{\Numer}[1]{N_{#1}}           % numeraire
\newcommand{\Fwd}[2]{F_{#1,#2}}          % forward price F_{t,T}
\newcommand{\Fwdpx}[2]{\mathrm{Fwd}_{#1}(#2)} % verbose forward-price F_t(T)

% --- Rates ---
\newcommand{\IRR}{\mathrm{IRR}}          % internal rate of return
\newcommand{\Yld}{y}                     % bond yield variable
\newcommand{\Strike}{K}                  % strike (price-space)
\newcommand{\Lock}{L}                    % locked yield / rate
\newcommand{\Repo}{r^{\mathrm{repo}}}    % repo rate
\newcommand{\Fund}{r^{\mathrm{fun}}}     % unsecured funding rate
\newcommand{\OIS}{r^{\mathrm{ois}}}      % OIS rate
\newcommand{\Coll}{r^{\mathrm{c}}}       % collateral rate
\newcommand{\Rcmt}{R^{\mathrm{cmt}}}     % constant-maturity-treasury rate
\newcommand{\Rswp}{R^{\mathrm{swp}}}     % swap rate (used for risky variant)
\newcommand{\Rec}{\mathrm{Rec}}          % recovery rate
\newcommand{\NPV}{\mathrm{NPV}}          % present-value functional
\newcommand{\PVone}{\mathrm{PV01}}       % risky annuity
\newcommand{\Loss}{\mathrm{Loss}}        % default-leg integral
\newcommand{\Sasw}{S^{\mathrm{ASW}}}     % asset-swap spread
\newcommand{\Scds}{S^{\mathrm{CDS}}}     % CDS spread
\newcommand{\Basis}{\mathrm{Basis}}      % CDS-bond basis

% --- Sensitivities ---
\newcommand{\Diff}[2]{D_{#1}^{#2}}       % Euler's notation: D_x^k[f]
\newcommand{\Riskf}{\mathrm{RiskFactor}} % bond risk factor (= -DV01*1e4)
\newcommand{\DVone}{\mathrm{DV01}}

% --- Indicator / misc ---
\newcommand{\Ind}{\mathbf{1}}
\newcommand{\given}{\,|\,}
\newcommand{\dd}{\,\mathrm{d}}

% --- Greeks-style ---
\newcommand{\Gam}{\Gamma}
\newcommand{\Vega}{\mathcal{V}}

% =====================================================================

% --- Measures and expectations ---
\newcommand{\Q}{\mathbb{Q}}              % risk-neutral measure
\newcommand{\Qt}{\mathbb{Q}^{T}}         % T-forward measure
\newcommand{\E}{\mathbb{E}}              % expectation
\newcommand{\Et}[1]{\E^{#1}}             % expectation under measure #1
\newcommand{\Ftime}[1]{\mathcal{F}_{#1}} % filtration at #1

% --- Discount and forward objects ---
\newcommand{\Disc}[2]{D_{#1,#2}}         % discount factor D_{t,T}
\newcommand{\Bond}[2]{P_{#1}(#2)}        % bond price P_t(y)
\newcommand{\Ann}[1]{A_{#1}}             % annuity / level
\newcommand{\Numer}[1]{N_{#1}}           % numeraire
\newcommand{\Fwd}[2]{F_{#1,#2}}          % forward price F_{t,T}
\newcommand{\Fwdpx}[2]{\mathrm{Fwd}_{#1}(#2)} % verbose forward-price F_t(T)

% --- Rates ---
\newcommand{\IRR}{\mathrm{IRR}}          % internal rate of return
\newcommand{\Yld}{y}                     % bond yield variable
\newcommand{\Strike}{K}                  % strike (price-space)
\newcommand{\Lock}{L}                    % locked yield / rate
\newcommand{\Repo}{r^{\mathrm{repo}}}    % repo rate
\newcommand{\Fund}{r^{\mathrm{fun}}}     % unsecured funding rate
\newcommand{\OIS}{r^{\mathrm{ois}}}      % OIS rate
\newcommand{\Coll}{r^{\mathrm{c}}}       % collateral rate
\newcommand{\Rcmt}{R^{\mathrm{cmt}}}     % constant-maturity-treasury rate
\newcommand{\Rswp}{R^{\mathrm{swp}}}     % swap rate (used for risky variant)
\newcommand{\Rec}{\mathrm{Rec}}          % recovery rate
\newcommand{\NPV}{\mathrm{NPV}}          % present-value functional
\newcommand{\PVone}{\mathrm{PV01}}       % risky annuity
\newcommand{\Loss}{\mathrm{Loss}}        % default-leg integral
\newcommand{\Sasw}{S^{\mathrm{ASW}}}     % asset-swap spread
\newcommand{\Scds}{S^{\mathrm{CDS}}}     % CDS spread
\newcommand{\Basis}{\mathrm{Basis}}      % CDS-bond basis

% --- Sensitivities ---
\newcommand{\Diff}[2]{D_{#1}^{#2}}       % Euler's notation: D_x^k[f]
\newcommand{\Riskf}{\mathrm{RiskFactor}} % bond risk factor (= -DV01*1e4)
\newcommand{\DVone}{\mathrm{DV01}}

% --- Indicator / misc ---
\newcommand{\Ind}{\mathbf{1}}
\newcommand{\given}{\,|\,}
\newcommand{\dd}{\,\mathrm{d}}

% --- Greeks-style ---
\newcommand{\Gam}{\Gamma}
\newcommand{\Vega}{\mathcal{V}}


\title{Note on: \emph{Why We Don't Discount at the ``Risk-Free Rate''}\\
       \large Anonymous expository note}
\author{Reading notes}
\date{\today}

\begin{document}
\maketitle

% =====================================================================
\section*{Metadata}
\begin{tabular}{@{}p{3.5cm}p{11cm}@{}}
\toprule
Title           & Why We Don't Discount at the ``Risk-Free Rate'' \\
Author(s)       & Not stated \\
Date            & Not stated \\
Source          & No SSRN/DOI; pedagogical note \\
Local file      & \texttt{papers/anon\_discounting\_textbooks/What\_Textbooks\_Get\_Wrong\_About\_Discounting\_1777005506.pdf} \\
Tags            & \texttt{discounting}, \texttt{repo}, \texttt{collateral},
                  \texttt{csa}, \texttt{ois}, \texttt{xva}, \texttt{multi-curve}, \texttt{expository} \\
Date read       & \today \\
\bottomrule
\end{tabular}

% =====================================================================
\section*{One-paragraph summary}
An expository piece arguing that the textbook ``discount at the
risk-free rate'' is a placeholder for two distinct rates: the repo
rate $\Repo$ of the hedge instrument, which enters the \emph{drift}
of the underlying under $\Q$ (and so the forward price); and the
collateral rate $\Coll$ specified in the CSA, which enters the
\emph{discount factor} applied to the derivative's value. The pre-2008
practice of collapsing both into a single LIBOR-flat rate was
operationally fine when LIBOR-OIS spreads were small; the crisis
forced an explicit split. The note works through three collateral
regimes (perfect CSA, no CSA, partial CSA), a worked 5-year USD swap
example showing a $\$570$k discount-rate impact on $\$100$m notional,
and a second channel by which repo enters discounting (non-cash
collateral). It is pedagogical rather than novel; equations exist
chiefly to anchor the worked examples.

% =====================================================================
\section{Setup and notation}

The note uses constant-rate examples throughout for pedagogical
clarity, while stating that production systems implement the
stochastic versions (eqs.~1--6 in the paper).

\begin{tabular}{@{}p{3cm}p{3.5cm}p{8.5cm}@{}}
\toprule
Paper            & House                  & Meaning \\
\midrule
$r(t)$           & ---                     & Short rate (textbook) \\
$r^{\text{repo}}(t)$ & $\Repo$             & Repo rate on the hedge instrument; drift under $\Q$ \\
$r_c(s)$         & $\Coll$                 & CSA / collateral rate; entered into discounting \\
$r_f(s)$         & $\Fund$                 & Unsecured funding rate of the dealer \\
$r_{\text{eff}}(s)$ & ---                  & Path-dependent effective rate under partial CSA \\
$V(t)$           & $V_t$                   & Derivative value \\
$S(t)$           & $S_t$                   & Underlying \\
$\sigma(t, S(t))$& $\sigma_t$              & Local volatility \\
$\lambda_C, \lambda_D$ & ---               & Default intensity of counterparty / dealer \\
$R_C, R_D$       & ---                     & Recovery rates of counterparty / dealer \\
$H$              & ---                     & CSA threshold \\
\bottomrule
\end{tabular}

% =====================================================================
\section{Key equations}

\paragraph{General stochastic pricing formula (uncollateralised).}
\begin{equation}\label{eq:general-pricing}
  V(t) \;=\; \Et{\Q}\!\left[ \exp\!\left(-\int_t^T r(s)\,\dd s\right) V(T) \,\Big|\, \Ftime{t} \right].
\end{equation}
\textit{Says:} the standard textbook formula --- but only when one
single rate $r(s)$ correctly captures both the cost of carry and the
discounting, which holds essentially never in modern practice.

\paragraph{Underlying dynamics under $\Q$ (with repo).}
\begin{equation}\label{eq:underlying-dyn}
  \dd S(t) \;=\; r^{\text{repo}}(t)\,S(t)\,\dd t \;+\; \sigma(t, S(t))\, S(t)\, \dd W^{\Q}(t).
\end{equation}
\textit{Says:} the drift of $S$ under the pricing measure is the
\emph{repo} rate of $S$, not a universal $r$. This is the cost of
carry of the replicating hedge.

\paragraph{Collateralised pricing.}
\begin{equation}\label{eq:coll-pricing}
  V(t) \;=\; \Et{\Q}\!\left[ \exp\!\left(-\int_t^T r_c(s)\,\dd s\right) V(T) \,\Big|\, \Ftime{t} \right].
\end{equation}
\textit{Says:} under a CSA with cash collateral earning $r_c$, the
discount rate is exactly $r_c$. This is forced by the no-arbitrage
martingale condition on $\exp(-\int_0^t r_c\,\dd s)\, V(t)$.

\paragraph{XVA components} (verbatim from paper):
\begin{align}
  \text{CVA} &= \Et{\Q}\!\left[\int_t^T e^{-\int_t^s r_c(u)\,\dd u}\,\lambda_C(s)\,(1-R_C)\,(V(s))^+ \dd s\,\Big|\,\Ftime{t}\right], \label{eq:cva}\\
  \text{DVA} &= \Et{\Q}\!\left[\int_t^T e^{-\int_t^s r_c(u)\,\dd u}\,\lambda_D(s)\,(1-R_D)\,(-V(s))^+ \dd s\,\Big|\,\Ftime{t}\right], \label{eq:dva}\\
  \text{FVA} &= \Et{\Q}\!\left[\int_t^T e^{-\int_t^s r_c(u)\,\dd u}\,(r_f(s) - r_c(s))\, V(s)\,\dd s\,\Big|\,\Ftime{t}\right]. \label{eq:fva}
\end{align}
\textit{Says:} CVA charges expected positive exposure for counterparty
default; DVA gives credit for the dealer's own default risk on negative
exposures; FVA charges the spread between funding and collateral rates
on the net exposure. All three discount at $r_c$, consistent with
\eqref{eq:coll-pricing}.

\paragraph{Pricing PDE with separated drift and discounting.}
For an equity derivative on $S$ with cash-collateralised PV:
\begin{equation}\label{eq:pde-split}
  \frac{\partial V}{\partial t}
  \;+\; r^{\text{repo}}\, S\, \frac{\partial V}{\partial S}
  \;+\; \tfrac{1}{2}\,\sigma^2\, S^2\, \frac{\partial^2 V}{\partial S^2}
  \;=\; r_c \cdot V.
\end{equation}
\textit{Says:} $r^{\text{repo}}$ in the drift term, $r_c$ on the RHS.
Collapsing them to a single $r$ (Black--Scholes) misses the post-2008
distinction.

\paragraph{Forward price under repo} (equity example):
\begin{equation}\label{eq:fwd-equity}
  F \;=\; S_0 \cdot e^{r^{\text{repo}} \cdot T}.
\end{equation}
\textit{Says:} the implied forward uses the repo rate, not the
risk-free rate. The example: $S_0 = \pounds 100$, $T = 1$,
$r^{\text{repo}} = 5.5\%$ gives $F \approx \pounds 105.65$, versus
the textbook $\pounds 105.13$ at $r = 5\%$ --- a $52$p difference
per share before any option valuation.

\paragraph{Three regimes for the discount rate}:
\begin{align}
  \text{Perfect CSA:} \quad V(t) &= e^{-r_c (T-t)}\, \Et{\Q}[V(T) | \Ftime{t}], \label{eq:reg-perfect}\\
  \text{No CSA:}      \quad V(t) &= e^{-r_f (T-t)}\, \Et{\Q}[V(T) | \Ftime{t}], \label{eq:reg-none}\\
  \text{Partial CSA:} \quad V(t) &= e^{-r_{\text{eff}} (T-t)}\, \Et{\Q}[V(T) | \Ftime{t}], \label{eq:reg-partial}
\end{align}
with $r_{\text{eff}}(s) = r_c(s)\,\Ind_{\{|V(s)| > H\}} + r_f(s)\,\Ind_{\{|V(s)| \le H\}}$.
\textit{Says:} the effective discount rate is path-dependent under
partial collateralisation --- which is why even constant-rate partial
CSAs require simulation.

\paragraph{Cross-currency CSA discrepancy.}
\begin{equation}\label{eq:xccy-csa}
  V^A(t) \;=\; e^{-r_c^{\text{EUR}} (T-t)}\, \Et{\Q}[V(T) | \Ftime{t}],
  \quad
  V^B(t) \;=\; e^{-r_c^{\text{USD}} (T-t)}\, \Et{\Q}[V(T) | \Ftime{t}].
\end{equation}
\textit{Says:} the same derivative has two different values under two
CSAs with different collateral currencies. The cross-currency basis
prices this difference.

\paragraph{Discount rate under non-cash (bond) collateral.}
\begin{equation}\label{eq:disc-bond-coll}
  r_{\text{eff}} \;=\; r^{\text{bond}}_{\text{repo}} + h,
\end{equation}
where $h$ is the haircut cost.
\textit{Says:} when the CSA permits bond collateral, the bank must
repo the bond to convert it into spendable cash; the effective
discount rate is the bond's repo rate (plus haircut), not the CSA
overnight rate. Distinct channel from the drift-side repo in
\eqref{eq:underlying-dyn}.

% =====================================================================
\section{Derivations \& commentary}

\begin{commentary}
The note's structural insight is the clean separation:
\emph{drift = repo, discount = collateral}. Both rates are typed
incorrectly in textbook treatments under the universal symbol $r$.
The summary table in \S3.3 is the operational decision rule and is
worth lifting into one's own notation:
\begin{center}
\begin{tabular}{@{}p{4cm}p{5cm}p{5cm}@{}}
\toprule
                  & Repo $\Repo$                    & Collateral $\Coll$ \\
\midrule
Role              & Drift of underlying under $\Q$  & Discount factor on $V$ \\
Determined by     & Repo market of the hedge asset  & CSA contract \\
Can differ?       & Yes, always                     & Yes, always \\
\bottomrule
\end{tabular}
\end{center}
\end{commentary}

\begin{commentary}
The two repo channels (\S3 drift and \S6.5 discount-via-non-cash-collateral)
are easy to conflate. They involve different bonds: the drift channel
uses the repo rate of \emph{the hedge instrument}; the discount channel
uses the repo rate of \emph{the bond posted as collateral}. These will
generally trade at different rates on the same day, especially when
specialness is involved.
\end{commentary}

\begin{commentary}
The note's \S6.3 worked example ($5$-year receiver swap, $\$100$m,
fixed rate $3.50\%$, $r_c = 4\%$ OIS, $r_f = 4.8\%$ unsecured) gives:
PV of fixed leg under perfect CSA $= \$15.54$m; no CSA $= \$14.97$m
($-\$570$k); partial CSA with $\pounds 10$m threshold $\approx \$15.29$m
($-\$250$k). Useful as a sanity-check anchor for any multi-curve
implementation.
\end{commentary}

% =====================================================================
\section{Validation}

\begin{validation}
This is a \emph{specification} of mathematical objects the paper
defines and the numerical values they should take. It says nothing
about how those objects are implemented or invoked.

\textbf{Quantities the paper defines:}
\begin{itemize}[noitemsep]
  \item Three regime-specific present-value formulas
        (eqs.~\ref{eq:reg-perfect}, \ref{eq:reg-none}, \ref{eq:reg-partial}).
  \item Equity forward price under repo (eq.~\ref{eq:fwd-equity}).
  \item Bond forward price under GC repo with coupon carry
        ($F^{\text{Gilt}} = P_0 \cdot e^{(r_{GC} - c) T}$, paper \S3.2).
  \item CVA, DVA, FVA integrals (eqs.~\ref{eq:cva}--\ref{eq:fva}).
  \item Effective discount rate under bond collateral
        (eq.~\ref{eq:disc-bond-coll}).
\end{itemize}

\textbf{Canonical numerical cases} (from worked examples in the paper):

\smallskip
\textit{Case A --- Equity forward, \S3.1:}
$S_0 = \pounds 100$, $T = 1$, $r^{\text{repo}} = 5.5\%$, $r = 5\%$ (textbook comparison).

\smallskip
\textit{Case B --- 5Y receiver swap, \S6.3:}
$\$100$m notional, fix rate $3.50\%$ received, SOFR paid,
$r_c = 4.00\%$, $r_f = 4.80\%$, $5$Y tenor.

\smallskip
\textit{Case C --- Bond collateral ColVA, \S6.5:}
$\pounds 20$m Gilt posted as collateral on a 10Y swap, duration $\approx 8.5$,
GC repo $4.95\%$ vs.\ SONIA $5.00\%$, or special repo $3.50\%$ vs.\
SONIA $5.00\%$.

\textbf{Expected numerical values:}

\smallskip
\begin{tabular}{@{}p{6cm}rl@{}}
\toprule
Quantity                                                & Value          & Tolerance \\
\midrule
Case A: equity forward $F$                              & $\pounds 105.65$ & $\pm 0.01$ \\
Case A: textbook (mis)forward                           & $\pounds 105.13$ & $\pm 0.01$ \\
Case B: PV fixed leg, perfect CSA                       & $\$15.54$m      & $\pm \$0.01$m \\
Case B: PV fixed leg, no CSA                            & $\$14.97$m      & $\pm \$0.01$m \\
Case B: impact vs.\ OIS, no CSA                         & $-\$570$k       & $\pm \$10$k \\
Case B: PV fixed leg, partial CSA (\$10m thresh.)       & $\$15.29$m      & $\pm \$0.01$m \\
Case C: ColVA, GC repo case                             & $\pounds 85$k    & order-of-mag. \\
Case C: ColVA, special repo case                        & $\pounds 2.55$m  & order-of-mag. \\
\bottomrule
\end{tabular}
\smallskip

\textbf{Equation $\leftrightarrow$ quantity map:}

\smallskip
\begin{tabular}{@{}p{2.8cm}p{3.8cm}p{6.2cm}@{}}
\toprule
Paper eq.                  & Symbol         & Quantity \\
\midrule
\eqref{eq:underlying-dyn}  & $\dd S(t)$     & dynamics of underlying with repo drift \\
\eqref{eq:coll-pricing}    & $V(t)$         & present value under perfect CSA \\
\eqref{eq:cva}--\eqref{eq:fva} & CVA, DVA, FVA & XVA components (all discount at $r_c$) \\
\eqref{eq:pde-split}       & ---            & pricing PDE with separated drift and discount \\
\eqref{eq:fwd-equity}      & $F$            & equity forward at repo rate \\
\eqref{eq:reg-perfect}--\eqref{eq:reg-partial} & $V(t)$ & PV in three collateral regimes \\
\eqref{eq:xccy-csa}        & $V^A, V^B$     & PVs under two different CSA currencies \\
\eqref{eq:disc-bond-coll}  & $r_{\text{eff}}$ & effective discount under bond collateral \\
\bottomrule
\end{tabular}
\smallskip

\textbf{Invariants and identities to verify:}
\begin{enumerate}[noitemsep]
  \item \emph{No-arbitrage on collateralised process}: the process
        $\exp(-\int_0^t r_c\,\dd s)\, V(t)$ is a $\Q$-martingale under
        a perfect CSA; equivalently, discounting at any rate
        $\tilde r \neq r_c$ produces a deterministic drift $(r_c - \tilde r) V(t)$,
        which is an arbitrage.
  \item \emph{Forward reconstruction (Case A)}: $S_0 e^{r^{\text{repo}} T}
        = \pounds 105.65$ at the canonical inputs within $\pm 0.01$.
  \item \emph{Discount-rate impact (Case B)}: changing the discount
        rate from $r_c = 4\%$ to $r_f = 4.8\%$ on a $\$100$m
        $5$-year swap fixed leg at $3.50\%$ moves the PV by
        $\approx -\$570$k.
  \item \emph{Limit consistency}: the partial-CSA formula
        \eqref{eq:reg-partial} reduces to \eqref{eq:reg-perfect}
        as $H \to 0$, and to \eqref{eq:reg-none} as $H \to \infty$.
  \item \emph{Cross-currency CSA disagreement}: \eqref{eq:xccy-csa}
        produces a strict inequality $V^A \neq V^B$ whenever
        $r_c^{\text{EUR}} \neq r_c^{\text{USD}}$, with the sign of
        the gap determined by the sign of $V(T)$ and the rate gap.
  \item \emph{Specialness $\to$ ColVA monotonicity (Case C)}: ColVA in
        eq.~\eqref{eq:disc-bond-coll} is decreasing in
        $r^{\text{bond}}_{\text{repo}}$ at fixed $r_c$. The
        $\pounds 85$k $\to \pounds 2.55$m jump between GC and special
        cases reflects a $145$~bp rate gap on $\pounds 20$m at
        duration $8.5$.
\end{enumerate}
\end{validation}

% =====================================================================
\section{Glossary of symbols (paper's notation)}
\begin{tabular}{@{}p{3cm}p{12cm}@{}}
\toprule
$r(t)$           & Generic short rate (textbook) \\
$r^{\text{repo}}(t)$ & Instantaneous repo rate process; cost of carry for the hedge \\
$r_c(s)$         & Stochastic collateral / CSA rate \\
$r_f(s)$         & Stochastic unsecured funding rate of the dealer \\
$r_{\text{eff}}(s)$ & Path-dependent effective rate under partial CSA \\
$V(t)$           & Derivative value \\
$S(t)$           & Underlying \\
$\sigma(t, S(t))$& Local volatility \\
$W^{\Q}(t)$      & Brownian motion under risk-neutral measure $\Q$ \\
$\lambda_C(t)$, $\lambda_D(t)$ & Default intensity of counterparty / dealer \\
$R_C, R_D$       & Recovery rates of counterparty / dealer \\
$H$              & CSA threshold (collateral floor) \\
$P_0, c$         & Bond clean price, coupon (in Gilt example) \\
$h$              & Haircut cost under non-cash collateral \\
SOFR, SONIA, €STR & Overnight rates in USD, GBP, EUR \\
GC               & General collateral (repo) \\
ColVA            & Collateral valuation adjustment \\
\bottomrule
\end{tabular}

% =====================================================================
\begin{todobox}
\begin{itemize}[noitemsep]
  \item Cross-check the Case B numbers ($\$15.54$m perfect CSA,
        $\$14.97$m uncollateralised) against a vanilla swap pricer
        with a hand-built OIS-only flat curve at $4\%$ and $4.8\%$.
  \item Implement a partial-CSA simulator that reproduces the
        $\$15.29$m blended PV for the \$10m-threshold case.
  \item Tab.~6.5 (Case C) gives a $\pounds 2.55$m ColVA gap purely
        from collateral specialness. Reproduce on a simple 10Y swap
        with bond-collateral option.
\end{itemize}
\end{todobox}

\end{document}
